\section{结论}
本文以\textbf{完整发现、全部清除和无遗漏退出}为先决条件，建立有界误差集合定位、保证接收选点和证书约束调度模型。问题一以可复现三角形反例说明直径不超过 $40$ m 并不足以保证清除；问题二的 $120$ 场景配对显示，稀疏候选平均节省 $5.74$ s，但平均后验半径增加 $3.68$ m，揭示局部精度与完整行动成本的权衡。

问题三以七站圆盘覆盖保证全向发现，问题四以 $31$ 站三角网格处理未知发射方向。最终模型在两问各 $3000$ 个独立研究场景中均全部完成，共清除 $38961$ 和 $39086$ 个源，\textbf{完整任务完成率和源级清除率均为 $100\%$}；含无遗漏确认的平均每源时间分别为 $255.99$ s 与 $767.25$ s。Q3 相对父模型的平均每源时间降低约 $0.453\%$；Q4 仅有小幅均值改善，配对区间跨零，未证明显著收益。

有利布局及规划对照表明，最后一个已发现源的清除时刻与完整搜索终点必须区分。$N<16$ 时未知频道仍需认证，发现数达到 $16$ 时才可利用数量上界停止其余覆盖。本文给出的最快构造样例用于解释这一完整性成本，并非所有策略的耗时下界。最终成绩仍须结合两问各三次正式测试记录核验。
